\section{Minimal Rotation}
\label{sec:cses-minimal-rotation}

\problemstatement{Given a string $s$, consider all $n$ rotations
(cyclic shifts) of it. Output the lexicographically smallest rotation.}

\begin{patternbox}
\emph{``Lexicographically smallest cyclic rotation''} has a linear-time
solution -- \textbf{Booth's algorithm} (or the equivalent least-rotation
scan) -- that finds the starting index of the minimal rotation by
comparing candidate starts and skipping ranges that cannot win.
\end{patternbox}

\subsection*{The least-rotation scan}

Work on $s+s$ (so every rotation is a length-$n$ window). Maintain two
candidate start positions $i$ and $j$ and an offset $k$: compare
$s[i+k]$ with $s[j+k]$. If they are equal, extend $k$. If one is smaller,
the losing candidate -- and, crucially, \emph{every start it skipped over}
-- can be discarded by jumping it past $j+k$ (or $i+k$). Because each
comparison either advances $k$ or eliminates a range of candidates, the
total work is linear. The surviving candidate is the minimal rotation's
start.

\begin{ideabox}[Skip Doomed Starts in Bulk]
When candidate $j$ beats candidate $i$ after matching $k$ characters,
none of the starts strictly between $i$ and $i+k+1$ can be minimal either
(they were dominated), so we jump $i$ past them all. This bulk elimination
is what keeps the scan $O(n)$ despite comparing rotations.
\end{ideabox}

\subsection*{Algorithm}

\begin{enumerate}
  \item Set $i=0$, $j=1$, $k=0$ over the doubled string $ss=s+s$.
  \item While all three are in range, compare $ss[i+k]$ and $ss[j+k]$:
        equal $\Rightarrow k{+}{+}$; otherwise advance the larger
        candidate past $\max(i,j)+k+1$, reset $k=0$, and keep $i<j$.
  \item The answer starts at $\min(i,j)$; output $ss[\text{start}\,..\,
        \text{start}+n-1]$.
\end{enumerate}

\subsection*{Dry run}

\dryrun{$s=$``acab''.}

\begin{itemize}
  \item Rotations: ``acab'',``caba'',``abac'',``baca''.
  \item Lexicographically smallest is ``abac'', starting at index $2$.
  \item The scan converges on start $=2$.
\end{itemize}

\subsection*{C++ implementation}

\begin{lstlisting}
#include <bits/stdc++.h>
using namespace std;

int leastRotation(const string& s) {
    string ss = s + s;
    int n = ss.size();
    int i = 0, j = 1, k = 0;
    while (i < (int)s.size() && j < (int)s.size() && k < (int)s.size()) {
        char a = ss[i + k], b = ss[j + k];
        if (a == b) { k++; }
        else if (a > b) { i = max(i + k + 1, j + 1); k = 0; }
        else            { j = max(j + k + 1, i + 1); k = 0; }
    }
    return min(i, j);                             // start of minimal rotation
}

int main() {
    string s;
    cin >> s;
    int start = leastRotation(s);
    string ss = s + s;
    cout << ss.substr(start, s.size()) << "\n";
    return 0;
}
\end{lstlisting}

\begin{complexitybox}
\textbf{Time Complexity.} $O(n)$ -- each comparison advances $k$ or a
candidate pointer, and both are bounded by $O(n)$ total. \textbf{Space
Complexity.} $O(n)$ for the doubled string.
\end{complexitybox}

\begin{takeawaybox}
The least-rotation scan finds the minimal cyclic shift in linear time by
eliminating dominated starts in bulk. It is the canonical tool for
canonicalising cyclic strings -- useful for necklace counting and
comparing circular sequences up to rotation.
\end{takeawaybox}
